\section{Constraints, active-set changes, and cutoff reuse}\label{sec:constraints}

Constraints can remove the directions responsible for a stall, but they
also create active-set changes that invalidate a local extrapolation.
Parametric LP duality gives precise contracts for some of this behavior.
These are uses of classical sensitivity analysis; the important point
here is the different conclusions supported by primal witnesses, dual
vectors, and complete basis regions.

\subsection{A fixed matrix and a changing right-hand side}

Consider the family
\begin{equation}\label{eq:parametric-lp}
 R(\theta)=\{z:Az\le b+E\theta\},\qquad
 h_c(\theta)=\max\{c^Tz:z\in R(\theta)\}.
\end{equation}
All parameters considered have nonempty relaxations and finite attained
support values. The original variables are selected coordinates of $z$;
signed coordinate objectives $c=\pm e_i$ give OBBT endpoints. Exact
certificates use rational data. Validity as an outer relaxation of the
original nonlinear model is a separate premise.

The matrix $A$ is fixed. Rebuilding McCormick rows generally changes their
coefficients, so~\eqref{eq:parametric-lp} does not automatically describe
box changes. It applies directly to cutoff changes with a frozen box.
Varying boxes require a valid normalization or separate verification of
the coefficient dependence. A numerical basis is only a proposal for
the checks below.

\begin{proposition}[Dual envelope]\label{prop:dual}
If $y\ge0$ and $A^Ty=c$, then
\begin{equation}\label{eq:dual-envelope}
 h_c(\theta)\le y^Tb+y^TE\theta.
\end{equation}
The minimum of finitely many such affine expressions is also an upper
bound, even when their original bases cease to be optimal or feasible.
If only the cutoff row $a^Tz\le U$ changes, a primal-dual pair is
exactly optimal at $U_0$, and $\eta\ge0$ is its cutoff multiplier, then
\begin{equation}\label{eq:cutoff-dual}
 h_c(U)\le h_c(U_0)-\eta(U_0-U)\qquad(U\le U_0).
\end{equation}
\end{proposition}
\begin{proof}
For feasible $z$, $c^Tz=y^TAz\le y^T(b+E\theta)$. Maximizing and
then taking the minimum over stored vectors proves the first statement.
For a single changing cutoff, the dual affine expression has slope
$\eta$ and agrees with $h_c(U_0)$ at the certified optimum, giving the
second statement.
\end{proof}

The quantity $\eta(U_0-U)$ is a guaranteed support reduction, not an
upper bound on possible tightening. A zero multiplier provides no such
guarantee; it does not prove that a larger cutoff change is ineffective.
Without exact optimality at $U_0$, the dual upper envelope remains valid,
but subtracting from an unverified numerical optimum does not establish
the displayed reduction.

\begin{proposition}[Exact basis region]\label{prop:basis}
Let $d=\dim z$, and choose $d$ linearly independent rows $I$. Define
\[
 z_I(\theta)=A_I^{-1}(b_I+E_I\theta),\qquad
 y_I=A_I^{-T}c,\qquad y_j=0\ (j\notin I).
\]
If $y_I\ge0$, then on the closed rational polyhedron
\[
 C_I=\{\theta:Az_I(\theta)\le b+E\theta\}
\]
the support is exactly
\begin{equation}\label{eq:basis-value}
 h_c(\theta)=c^TA_I^{-1}(b_I+E_I\theta).
\end{equation}
\end{proposition}
\begin{proof}
The point $z_I(\theta)$ is feasible on $C_I$. The vector $y$ is dual
feasible and supported on rows tight at that point, so the primal and
dual objective values agree. Proposition~\ref{prop:dual} supplies the
matching upper bound.
\end{proof}

Zero multipliers and ties are allowed; singular proposed bases are
rejected. Checking the affine residuals at every vertex of a proposed
parameter polytope proves its containment in $C_I$. This does not prove
that a collection of such polytopes covers all parameters of interest.
For a scalar cutoff, substituting $z_I(U)=v+wU$ makes every residual an
affine inequality in $U$; their exact intersection is a validity interval.
Beyond it, the stored dual expression is still an upper bound but need
not equal the support.

For an exact example, minimize $y$ subject to $0\le x\le3$,
$0\le y\le4$, and
\[
 x-y\le1,\qquad 4x-y\le7,\qquad x\le5/2.
\]
With cutoff $y\le U$, $0\le U\le4$, the upper support of $x$ is
\begin{equation}\label{eq:cutoff-example}
 h_x(U)=\min\{1+U,7/4+U/4,5/2\}
 =\begin{cases}
 1+U,&0\le U\le1,\\
 7/4+U/4,&1\le U\le3,\\
 5/2,&3\le U\le4.
 \end{cases}
\end{equation}
Choosing $y=U$ and the displayed $x$ gives a feasible point meeting each
upper bound, proving the formula. At $U_0=4$ the cutoff multiplier is
zero, yet moving to $U=2$ reduces the support from $5/2$ to $9/4$.
The witness $(5/2,3)$ proves no reduction only down to cutoff $3$.
At $U_0=2$, extrapolating the old dual to $U=1/2$ gives the safe upper
bound $15/8$, while the actual support is $3/2$.

\subsection{Combining retained primal and dual information}

For a signed coordinate objective $c_i$, let $p_i$ be the corresponding
current box endpoint. Keep all rows except the cutoff fixed. From stored
feasible relaxed points $W_i$, retain those satisfying the new cutoff and
define, when that subset is nonempty,
\[
 L_i(U)=\max\{c_i^Tz:z\in W_i,\ a^Tz\le U\}.
\]
If $F_i(U)$ is the exact support and $D_i(U)$ is a stored dual envelope,
then
\begin{equation}\label{eq:bracket}
 L_i(U)\le F_i(U)\le\min\{p_i,D_i(U)\},\qquad
 0\le p_i-F_i(U)\le p_i-L_i(U).
\end{equation}
The left bound follows from primal feasibility, the right from duality
and the current box. A near-endpoint witness certifies little possible
one-round improvement; a sufficiently low dual envelope certifies a
useful available improvement. The interval left between them is a natural
place for a budgeted probe.

An old support-attaining witness proves the endpoint unchanged until its
objective value is excluded. Finding the lowest-objective witness among
old optima can itself require another LP, whose cost must be charged.
An excluded witness does not prove that the support moves: another
witness may remain. Every changed local row requires rechecking, and
screening a whole round requires information for every direction whose
benefit is bounded.

\subsection{A uniform bound across active-set changes}

\begin{proposition}[Complete basis coverage implies a comparison matrix]\label{prop:cells}
Fix the cutoff. Let $\mathcal D$ be a convex endpoint region for which,
for each support coordinate $i$, finitely many verified closed basis
regions cover $\mathcal D$. Suppose on each region
$F_i(p)=\alpha_i+g_i^Tp$. If $M\ge0$ and $|g_{ij}|\le M_{ij}$
for every covered region, then
\[
 |F(p)-F(q)|\le M|p-q|\qquad(p,q\in\mathcal D).
\]
\end{proposition}
\begin{proof}
The segment from $p$ to $q$ stays in $\mathcal D$. Its intersections
with the finitely many polyhedral regions partition it into finitely
many intervals. On each, the derivative of coordinate $i$ along the
segment has magnitude at most $\sum_jM_{ij}|p_j-q_j|$. At boundaries,
both basis formulas equal the same optimal support, so there is no jump.
Integrating over the segment gives the claim.
\end{proof}

This provides the uniform premise in Theorem~\ref{thm:residual}. It
does not prove invariance of $\mathcal D$ or existence of a finite
majorant $e$. For the order interval between a protected fixed box and
the current box, monotonicity supplies invariance. A derivative bound
only on the present basis region is insufficient if later rounds leave
that region.

A deliberately simple example verifies all these steps. For $x^2\le0$,
$0\le x\le r$, retain only tangents to $x^2$ at
$t_1=r/2$ and $t_2=r/8+3/4$. For positive tangent location $t$, its
necessary inequality $2tx-t^2\le0$ is $x\le t/2$. Together with the
box, the upper-support map on $0\le r\le8$ is
\begin{equation}\label{eq:switch-map}
 F(r)=\min\{r/4,r/16+3/8\}.
\end{equation}
At $r=0$ the box fixes $x=0$, so no division by a zero tangent
coefficient is needed. The two basis regions are $[0,2]$ and $[2,8]$,
with slopes $1/4$ and $1/16$. They cover the invariant interval, giving
the uniform bound $M=[1/4]$. The trajectory
\[
 8\longmapsto7/8\longmapsto7/32\longmapsto7/128\longmapsto\cdots
\]
crosses a basis boundary. Its first width ratio $7/64$ is smaller than
the subsequent ratio $1/4$. The example tests a sensitivity certificate;
direct nonlinear propagation would solve the original problem at once.

\subsection{An equality changes the exact contraction rate}

Consider $f(x,y)=(x-y)^2$, expanded as $x^2+y^2-2xy$, on $[-r,r]^2$.
Keep the convex square terms exact and use the McCormick overestimator
for $xy$. The relaxed objective is
\[
 \phi_r(x,y)=x^2+y^2-2r^2+2r|x-y|.
\]
Without further constraints the minimizers $(r,r)$ and $(-r,-r)$
preserve every endpoint at cutoff zero, so OBBT stalls. Impose instead
the exact equality $y=-x$. On it,
$\phi_r(t,-t)=2t^2+4r|t|-2r^2$.

\begin{proposition}[Constrained exact rate and cutoff floor]\label{prop:equality}
For cutoff $U\ge0$, a Jacobi round on the symmetric box returns radius
\[
 G_U(r)=\min\{r,\sqrt{2r^2+U/2}-r\}.
\]
At $U=0$, the exact linear rate is $\sqrt2-1$. At $U>0$, from any
$r_0\ge\sqrt U/2$, the radii decrease to exactly $\sqrt U/2$.
\end{proposition}
\begin{proof}
Solving $2|t|^2+4r|t|-2r^2\le U$ gives the stated root, restricted
by the current radius. Set $a=\sqrt U/2$. The root is less than $r$
precisely for $r>a$, and at least $a$ because
$2r^2+2a^2-(r+a)^2=(r-a)^2\ge0$. Hence the radii decrease to a
limit at least $a$; continuity and the fixed-point equation force that
limit to be $a$. Factoring $r$ gives the zero-cutoff rate.
\end{proof}

Eliminating the equality before relaxing gives the objective $4x^2$.
Its exact convex cutoff reaches $|x|\le\sqrt U/2$ in one round.
Thus the rate depends on the formulation and relaxation, not only on
the original constrained problem. Affine equalities can also restrict
the earlier tangent support problems to a nullspace when the original
expansion has no first-order objective term and is uniform on the
admitted shapes. General nonlinear active constraints additionally need
uniform residual and remainder control. A first-order tangent cone
alone does not establish a finite contraction theorem.

\subsection{Nonlinear constraints through a feasible repair}

A uniform constraint error bound provides a route that does not require
guessing the next active set. Let $S$ be the original feasible set and
$x^*\in S$ a minimizer with value $f^*$. Consider boxes containing
$x^*$ in a fixed neighborhood, and a nonnegative constraint residual
$v(x)$ that vanishes on $S$. Suppose every point $x$ in each projected
relaxation satisfies
\begin{equation}\label{eq:repair-data}
 \phi_B(x)\ge f(x)-a w(B)^2,\qquad v(x)\le b w(B)^2,
\end{equation}
with uniform $a,b\ge0$. Assume it has a feasible repair $y\in S$ such
that, for uniform $\kappa,L\ge0$, $0<\theta\le1$, and $\mu>0$,
\begin{align}
 \norm{x-y}_\infty&\le\kappa v(x)^\theta,\label{eq:repair-distance}\\
 f(y)&\le f(x)+L\norm{x-y}_\infty,\label{eq:repair-lipschitz}\\
 f(y)&\ge f^*+\mu\norm{y-x^*}_\infty^2.\label{eq:repair-growth}
\end{align}
The repair need not stay in $B$, but must stay in the neighborhood where
the last two estimates hold. These are uniform assumptions on all relaxed
points, including infeasible ones, not consequences of a numerical KKT
residual. The argument is a standard transfer of error and growth bounds;
an antecedent in the repository is
\href{run:../../notes/research-20260922-error-bound-transfer.md}{the error-bound transfer note}.

\begin{theorem}[Constrained width bound from a feasible repair]\label{thm:repair}
Under~\eqref{eq:repair-data}--\eqref{eq:repair-growth}, with
$U=f^*+\epsilon$, $\epsilon\ge0$, set
$r=\kappa(bw(B)^2)^\theta$. Then
\begin{equation}\label{eq:repair-width}
 w(T_U(B))\le2r+2\sqrt{\frac{\epsilon+aw(B)^2+Lr}{\mu}}.
\end{equation}
For $\theta=1$, let $K=\kappa b$ and
$q=2\sqrt{(a+LK)/\mu}$. If $q<1$ and
$\lambda=q+2K\overline w<1$ for some $\overline w>0$, every such
box with $w(B)\le\overline w$ satisfies
\begin{equation}\label{eq:repair-recurrence}
 w(T_U(B))\le\lambda w(B)+2\sqrt{\epsilon/\mu}.
\end{equation}
For nested iterates in the stated neighborhood,
\[
 \limsup_k w(B_k)\le\frac{2\sqrt{\epsilon/\mu}}{1-\lambda}.
\]
At zero cutoff slack ($\epsilon=0$), the widths contract with upper
ratio $\lambda$. If $K=0$, take $\lambda=q$ on any admitted width
neighborhood.
\end{theorem}
\begin{proof}
For a retained relaxed point, $f(x)\le f^*+\epsilon+aw(B)^2$.
Its repair has distance at most $r$ and objective at most
$f^*+\epsilon+aw(B)^2+Lr$. Feasible growth bounds its distance to
$x^*$ by the square root in~\eqref{eq:repair-width}. The triangle
inequality bounds the distance of $x$ by this radius plus $r$.
Taking coordinate ranges doubles the bound. If $\theta=1$, split the
square root using $\sqrt{s+t}\le\sqrt s+\sqrt t$ and bound
$2Kw(B)^2\le2K\overline w\,w(B)$. Nesting keeps all later widths
at most $\overline w$, and summing the scalar geometric recurrence
gives the stated limit bound.
\end{proof}

The limit bound is an upper bound, not an exact floor. At zero cutoff
slack, with
$\theta<1$ and fixed positive $L\kappa b$, the square-root term
in~\eqref{eq:repair-width} has order $w^\theta$, so this argument
need not establish linear contraction. This is a limitation of the
estimate, not proof of a stall. A local repair-pair constant $L_B$ of
order $w^{2-2\theta}$ can recover an $O(w)$ square-root term, but
the separate repair term $O(w^{2\theta})$ also needs control.

An explicit nonlinear equality illustrates the uniform assumptions.
Let
\[
 S=\{(x,y):y=x^2\}\cap([-1/4,1/4]\times[0,1/16]),
 \qquad f(x,y)=x^2+y^2-xy/16.
\]
The minimizer is $(0,0)$ with value zero. For a subbox
$B=[\ell,u]\times[0,v]$ containing the origin, use
\[
 x^2\le y\le(\ell+u)x-\ell u,\qquad
 \phi_B(x,y)=x^2+y^2-\frac1{16}
              \min\{uy,\ell y+vx-\ell v\}.
\]
The square epigraph and secant relax the graph, and the displayed
bilinear overestimator gives a convex objective estimator. A relaxed
point has feasible repair $(x,x^2)$ in the same box, with distance
$y-x^2\le(u-\ell)^2/4\le w(B)^2/4$.
The bilinear-envelope error is at most $(u-\ell)v/4$, so one may take
\[
 a=1/64,\quad b=1/4,\quad\kappa=1,\quad\theta=1.
\]
On the outer box,
$|\partial_x f|+|\partial_y f|\le(2+1/16)(1/4+1/16)=165/256$,
so take $L=165/256$. On the feasible graph,
\[
 f(x,x^2)=x^2(1-x/16+x^2)\ge(63/64)x^2
        =(63/64)\norm{(x,x^2)}_\infty^2.
\]
Thus $\mu=63/64$, $K=1/4$, and
$q^2=181/252<49/64$. For $w(B)\le1/8$, the rational upper rate
$\lambda=15/16$ gives
\[
 w(T_U(B))\le\frac{15}{16}w(B)+2\sqrt{64\epsilon/63}.
\]
This covers every admitted nested box shape without choosing an active
set. Its constants are conservative. The example uses exact convex
square epigraphs, which the production finite-tangent LP does not supply;
it proves a constrained theorem, not the rate of the measured plugin.

\subsection{Finite observed histories do not certify an eventual stall}

Fix $N\ge2$, and set $a=1/2$, $\delta=1/(2N)$, $b=a-\delta$,
and $q=1/4$. On $[0,1]$ define
\[
 F_{\rm stall}(s)=\min\{s,\max\{b,s-\delta\}\},
\]
and
\[
 F_{\rm shrink}(s)=
 \begin{cases}
 qs,&s\le b,\\
 qb+(b-qb)(s-b)/(a-b),&b\le s\le a,\\
 s-\delta,&s\ge a.
 \end{cases}
\]
Both maps are continuous and nondecreasing, with $0\le F(s)\le s$;
these assertions follow from the nonnegative slopes and the endpoints
of each affine piece. Thus $R_s=[0,F(s)]$ is a valid monotone outer
relaxation family for the feasible set $\{0\}$.
Starting at $s_0=1$, both trajectories subtract $\delta$ through
$s_{N+1}=b$. One then stalls at $b$, and the other converges
geometrically to zero. By increasing $N$, arbitrarily many common early
ratios can be made close to one despite the different eventual outcomes.

These are not asserted to be standard McCormick relaxations. They show
that validity, monotonicity, and a finite observed history alone cannot
certify the remaining benefit. A protected box, a uniform sensitivity
bound, or further relaxation structure supplies information absent from
that history. Exact checks of the examples and basis procedures appear
in \pathref{../theory/check_constrained_obbt.py}{theory/check_constrained_obbt.py}.
